\section{Conclusion}
\label{sec:conclusion}

This project used the AGJK marked-vertex search algorithm as the QIS core of a
process-level model of cognitive search. We implemented the Szegedy and
interpolated walks, verified the $\widetilde O(\sqrt{\HT})$ scaling on four
graph families, and compared single, independent, and clustered targets. The
main structural finding is that independent targets are easier, while clustered
targets give only a modest reduction relative to a single target ($\approx3\%$
on the barbell, $\approx17\%$ on the $8\times8$ grid); the large $\HT^+\gg\HT$ gap requires the
adversarial periodic arrangement of AGJK rather than clustering alone. On the
cognitive side, a toy semantic graph produced distance-dependent activation
patterns, a multi-branch decision graph produced a robust similarity-like
context effect under option coupling that the classical walk lacks, and a partial-measurement model showed
that, beyond the shortest deadlines, an intermediate measurement rate that can
stop the search finds clustered targets faster than both the pure quantum and
the classical walk.

The broader lesson is that the algorithmic structure of quantum walk search is
a useful source of cognitive hypotheses. The marked-set structure, the
extended hitting time, and the success curve are not merely complexity
quantities; they define a process model with testable predictions. At the same
time, the project is deliberately cautious: the results are numerical,
model-level, and do not claim a biological quantum advantage. The next step is
to test the same predictions on larger semantic networks and, eventually, on
behavioral choice and response-time data.
